\chapter*{Conclusion}
\addcontentsline{toc}{chapter}{Conclusion}
\markboth{Conclusion}{}
% BUDGET : 1 page (condense le 2026-09-01, cf. memoire/ETAT.md).
% REGLE : la conclusion repond a la question de l'introduction DANS SES TERMES. Si les mots
% changent, c'est qu'on repond a une autre question.

Le modèle porté dépasse les trois métriques publiées par \textcite{chopin2015}.

\begin{itemize}[leftmargin=1.2em,itemsep=3pt]
  \item \SI{\calibRetenuTypes}{\percent} de types d'exploitation reproduits contre
    \SI{\mesArtTypes}{\percent}.
  \item \SI{\calibRetenuParcelles}{\percent} de parcelles contre \SI{\mesArtParcelles}{\percent}.
  \item \SI{\calibRetenuSurface}{\percent} de surface contre \SI{\mesArtSurface}{\percent}.
\end{itemize}

La canne à sucre est simulée à \SI{\simRetenuCS}{\ha} contre \SI{\obsHaCS}{\ha} observés,
soit \SI{\padRetenuCS}{\percent} d'écart. Aucune contrainte ne la nomme et elle occupe plus de
la moitié du territoire. Ce phénomène s'est également produit pour la banane revenant sous le
plafond de plantain. C'est un marqueur de qualité du modèle.
En revanche, \num{\calibRetenuCulturesOk} cultures sur \num{\calibRetenuCulturesTotal}
passent sous le seuil de l'article, contre \num{8} sur \num{10} publiées.
Le domaine de validité est connu. Le modèle gère beaucoup mieux les cultures massives que les
cultures marginales.

Sauf pour la canne et la banane, le modèle n'intègre pas la notion de débouché. C'est une
réserve importante.
La marge brute varie en sens inverse du plafond de subventions. On établit que l'aide publique
achète de la stabilité et pas la marge. Quand les subventions se resserrent, la banane
d'exportation a tendance à disparaître et la surface de canne à fortement diminuer au profit de
l'élevage et de cultures vivrières comme le maraîchage.
Le classement des politiques dépend ensuite du forçage. Une politique n'est pas bonne
dans l'absolu. Ce qui est intéressant est de la confronter à différents contextes pour évaluer
sa robustesse et sa performance.

Il existe plusieurs pistes d'améliorations du projet.

\begin{itemize}[leftmargin=1.2em,itemsep=3pt]
  \item Modéliser les débouchés pour arrêter de supposer un marché parfaitement élastique à la
    hausse du vivrier.
  \item Inclure l'élevage pour mieux modéliser les prairies et affiner les indicateurs
    nutritionnels.
  \item Simuler les rendements plutôt que les tabuler pour lever l'écart du
    §~\ref{sec:rendements}. Cela est prévu dans la suite du stage à la suite d'une formation sur
    l'outil MAELIA de l'\gls{inrae}.
  \item Rendre le modèle dynamique afin de simuler les roulements de cultures d'une année à
    l'autre.
  \item Organiser des ateliers qui rassemblent les acteurs afin de concevoir les scénarios à
    explorer et critiquer leurs résultats.
\end{itemize}
